\chapter{Solutions to Chapter 28: Controlled Direct Effect}
\label{sol:chapter28}

\begin{exercisesolution}{$\CDE$ and $\NDE$}{cross-world independence 下 CDE 对 NDE 的加权表示}
本题的要点是：controlled direct effect 固定 mediator 的取值，而 natural direct effect 让 mediator 取它在 control treatment 下自然会取得的值。因此，若要把 $\NDE(x)$ 写成一组 $\CDE(m\mid x)$ 的平均，必须能把 $Y(z,m)$ 与 $M_0$ 在给定 $X=x$ 后分开。这正是题目给出的 cross-world independence 的作用。

固定 $X=x$。由于 $M$ 是离散变量，由全期望公式，
\[
E\{Y(z,M_0)\mid X=x\}
=
\sum_m E\{Y(z,m)\mid M_0=m,X=x\}\mathbb{P}(M_0=m\mid X=x).
\]
题设 cross-world independence 对所有 $z,z'$ 与 $m$ 都成立。取 $z'=0$，得到
\[
Y(z,m)\ind M(0)\mid X.
\]
因此
\[
E\{Y(z,m)\mid M_0=m,X=x\}
=
E\{Y(z,m)\mid X=x\}.
\]
于是
\[
E\{Y(z,M_0)\mid X=x\}
=
\sum_m E\{Y(z,m)\mid X=x\}\mathbb{P}(M_0=m\mid X=x).
\]
分别令 $z=1$ 与 $z=0$ 并相减，得到
\begin{align*}
\NDE(x)
&=
E\{Y(1,M_0)-Y(0,M_0)\mid X=x\}\\
&=
\sum_m \bigl[
E\{Y(1,m)\mid X=x\}
-
E\{Y(0,m)\mid X=x\}
\bigr]\mathbb{P}(M_0=m\mid X=x)\\
&=
\sum_m \CDE(m\mid x)\mathbb{P}(M_0=m\mid X=x).
\end{align*}
这就证明了题中的表示式。

若没有 cross-world independence，上述关系一般不成立。原因是 $\NDE(x)$ 中的权重不是简单地对固定的 response surface $Y(z,m)$ 做边际平均，而是对那些满足 $M_0=m$ 的 units 做条件平均：
\[
\NDE(x)
=
\sum_m
E\{Y(1,m)-Y(0,m)\mid M_0=m,X=x\}
\mathbb{P}(M_0=m\mid X=x).
\]
只有当
\[
E\{Y(1,m)-Y(0,m)\mid M_0=m,X=x\}
=
E\{Y(1,m)-Y(0,m)\mid X=x\}
\]
时，它才退化为题中的 $\CDE(m\mid x)$ 加权平均。没有 cross-world independence 时，$M_0$ 可以选择性地标记出具有不同 direct-effect response surface 的 units，因而这个等式没有理由成立。

\paragraph{Remark.}
这道题把 \cref{chapter::mediation,chapter::cde} 的差别压缩成一个公式：$\NDE$ 是以自然 mediator distribution 加权的 direct-effect contrast；$\CDE$ 是把 mediator 外力固定后的 direct-effect contrast。二者能否相连，关键不在代数，而在 cross-world independence 是否允许我们把 $M_0$ 与 $Y(z,m)$ 解耦。
\end{exercisesolution}

\begin{exercisesolution}{Observational studies with a multi-valued treatment}{多值 treatment 下 outcome mean 的识别与 doubly robust 公式}
本题证明 \cref{thm::treatment-multi-valued-Z}。固定一个 treatment level $k$。在 \cref{assume::ignorability-multi-valued-Z} 下，
\[
Z\ind \{Y(1),\ldots,Y(K)\}\mid X,
\]
特别地，
\[
Z\ind Y(k)\mid X.
\]
再由 consistency，当 $Z=k$ 时 observed outcome 满足 $Y=Y(k)$。

先证明 outcome regression identification。由全期望公式和 conditional ignorability，
\[
\mu_k
=
E\{Y(k)\}
=
E\bigl[E\{Y(k)\mid X\}\bigr].
\]
另一方面，
\begin{align*}
E(Y\mid Z=k,X)
&=
E\{Y(k)\mid Z=k,X\}\\
&=
E\{Y(k)\mid X\}.
\end{align*}
因此 $\mu_k(X)=E(Y\mid Z=k,X)$ 正好等于 $E\{Y(k)\mid X\}$，从而
\[
\mu_k=E\{\mu_k(X)\}.
\]

再证明 IPW identification。由于 overlap 保证 $e_k(X)>0$，有
\begin{align*}
E\left\{\frac{\mathbb{I}(Z=k)Y}{e_k(X)}\right\}
&=
E\left[
E\left\{\frac{\mathbb{I}(Z=k)Y}{e_k(X)}\mid X\right\}
\right]\\
&=
E\left[
\frac{1}{e_k(X)}
E\{\mathbb{I}(Z=k)Y(k)\mid X\}
\right]\\
&=
E\left[
\frac{1}{e_k(X)}
E\{\mathbb{I}(Z=k)\mid X\}
E\{Y(k)\mid X\}
\right]\\
&=
E\{Y(k)\}
=
\mu_k.
\end{align*}
第三行使用了 $Z\ind Y(k)\mid X$。

最后证明 doubly robust 性质。记
\[
m_k(X)=\mu_k(X,\beta),
\qquad
g_k(X)=e_k(X,\alpha).
\]
题中的 working estimand 为
\[
\mu_k^{\textup{dr}}
=
E\{m_k(X)\}
+
E\left[
\frac{\mathbb{I}(Z=k)\{Y-m_k(X)\}}{g_k(X)}
\right].
\]
若 outcome model 正确，即 $m_k(X)=\mu_k(X)=E\{Y(k)\mid X\}$，则
\begin{align*}
E\left[
\frac{\mathbb{I}(Z=k)\{Y-m_k(X)\}}{g_k(X)}
\mid X
\right]
&=
\frac{1}{g_k(X)}
E\{\mathbb{I}(Z=k)(Y(k)-m_k(X))\mid X\}\\
&=
\frac{e_k(X)}{g_k(X)}
\{E(Y(k)\mid X)-m_k(X)\}\\
&=0.
\end{align*}
因此 $\mu_k^{\textup{dr}}=E\{m_k(X)\}=E\{\mu_k(X)\}=\mu_k$。

若 propensity score model 正确，即 $g_k(X)=e_k(X)$，则
\begin{align*}
\mu_k^{\textup{dr}}
&=
E\{m_k(X)\}
+
E\left[
\frac{\mathbb{I}(Z=k)\{Y(k)-m_k(X)\}}{e_k(X)}
\right]\\
&=
E\{m_k(X)\}
+
E\left[
E\left\{
\frac{\mathbb{I}(Z=k)\{Y(k)-m_k(X)\}}{e_k(X)}
\mid X
\right\}
\right]\\
&=
E\{m_k(X)\}
+
E\{E(Y(k)\mid X)-m_k(X)\}\\
&=
E\{Y(k)\}
=
\mu_k.
\end{align*}
所以只要 treatment model 或 outcome model 二者之一正确，doubly robust 公式就等于目标均值。

\paragraph{Remark.}
\cref{thm::cde-identification} 只是把 joint treatment level 写成 $(Z,M)=(z,m)$。此时 $K=4$，$e_k(X)$ 变成 $e_{zm}(X)=\mathbb{P}(Z=z,M=m\mid X)$，而 $\CDE(m)$ 是两个 joint levels $(1,m)$ 与 $(0,m)$ 的 potential outcome means 之差。
\end{exercisesolution}

\begin{exercisesolution}{$\CDE$ in the linear outcome model}{线性 outcome model 中 interaction 决定 CDE 是否随 mediator 改变}
本题直接使用 \cref{assume::si-cde} 与 \cref{thm::cde-identification} 的 outcome-regression 表示。Sequential ignorability 给出
\[
E\{Y(z,m)\mid X\}=E(Y\mid Z=z,M=m,X).
\]
因此
\[
\CDE(m)
=
E\{Y(1,m)-Y(0,m)\}
=
E\bigl[
E(Y\mid Z=1,M=m,X)-E(Y\mid Z=0,M=m,X)
\bigr].
\]

若
\[
E(Y\mid Z,M,X)
=
\theta_0+\theta_1 Z+\theta_2 M+\theta_4\tran X,
\]
则
\[
E(Y\mid Z=1,M=m,X)
=
\theta_0+\theta_1+\theta_2 m+\theta_4\tran X,
\]
而
\[
E(Y\mid Z=0,M=m,X)
=
\theta_0+\theta_2 m+\theta_4\tran X.
\]
二者相减为 $\theta_1$，与 $X$ 和 $m$ 都无关。于是
\[
\CDE(m)=E(\theta_1)=\theta_1.
\]

若进一步允许 treatment-mediator interaction，
\[
E(Y\mid Z,M,X)
=
\theta_0+\theta_1 Z+\theta_2 M+\theta_3 ZM+\theta_4\tran X,
\]
则
\[
E(Y\mid Z=1,M=m,X)
=
\theta_0+\theta_1+\theta_2 m+\theta_3 m+\theta_4\tran X,
\]
而
\[
E(Y\mid Z=0,M=m,X)
=
\theta_0+\theta_2 m+\theta_4\tran X.
\]
相减得到
\[
\theta_1+\theta_3 m.
\]
因此
\[
\CDE(m)=\theta_1+\theta_3 m.
\]

\paragraph{Remark.}
这解释了 \cref{eg::cde-linear} 中的现象。没有 $ZM$ interaction 时，固定 mediator 的水平不会改变 treatment 的 additive contrast；一旦加入 interaction，direct effect 的大小随被固定的 mediator level 改变。
\end{exercisesolution}

\begin{exercisesolution}{$\CDE$ in the logistic outcome model}{logistic outcome model 中 CDE 是风险差而不是 log-odds coefficient}
本题与上一题形式相同，但二元 outcome 下 $\CDE(m)$ 仍然定义在 outcome mean scale 上：
\[
\CDE(m)=E\{Y(1,m)-Y(0,m)\}.
\]
当 $Y$ 是 binary outcome 时，$E\{Y(z,m)\mid X\}=\mathbb{P}\{Y(z,m)=1\mid X\}$。由 \cref{assume::si-cde}，
\[
\mathbb{P}\{Y(z,m)=1\mid X\}
=
\mathbb{P}(Y=1\mid Z=z,M=m,X).
\]

先考虑没有 interaction 的 logistic model：
\[
\logit\{\mathbb{P}(Y=1\mid Z,M,X)\}
=
\theta_0+\theta_1 Z+\theta_2 M+\theta_4\tran X.
\]
对两边取 inverse logit，得到
\[
\mathbb{P}(Y=1\mid Z=z,M=m,X)
=
\expit(\theta_0+\theta_1 z+\theta_2 m+\theta_4\tran X).
\]
因此
\begin{align*}
\CDE(m)
&=
E\bigl[
\mathbb{P}(Y=1\mid Z=1,M=m,X)
-
\mathbb{P}(Y=1\mid Z=0,M=m,X)
\bigr]\\
&=
E\{
\expit(\theta_0+\theta_1+\theta_2 m+\theta_4\tran X)
-
\expit(\theta_0+\theta_2 m+\theta_4\tran X)
\}.
\end{align*}
这就是题中第一式。

再考虑含 interaction 的 logistic model：
\[
\logit\{\mathbb{P}(Y=1\mid Z,M,X)\}
=
\theta_0+\theta_1 Z+\theta_2 M+\theta_3 ZM+\theta_4\tran X.
\]
同理，
\[
\mathbb{P}(Y=1\mid Z=z,M=m,X)
=
\expit(\theta_0+\theta_1 z+\theta_2 m+\theta_3 zm+\theta_4\tran X).
\]
于是
\begin{align*}
\CDE(m)
&=
E\{
\expit(\theta_0+\theta_1+\theta_2 m+\theta_3 m+\theta_4\tran X)\\
&\qquad\qquad
-
\expit(\theta_0+\theta_2 m+\theta_4\tran X)
\}.
\end{align*}
这正是题中第二式。

\paragraph{Remark.}
线性模型里 $\theta_1$ 本身就是 additive mean contrast；logistic 模型里 $\theta_1$ 是 conditional log-odds ratio 的一部分，而 $\CDE(m)$ 是 marginal risk difference。因此即使没有 $ZM$ interaction，$\CDE(m)$ 也通常依赖于 $m$ 和 $X$ 的分布。这个非线性差异是 logistic mediation 或 marginal structural model 中最容易混淆的地方。
\end{exercisesolution}
